% ==================== 1 晶体学与电子衍射基础 ====================
\section{晶体学与电子衍射基础}

\subsection{晶胞参数、晶胞体积与倒易点阵}

描述一个晶胞需要六个参数：三条棱长 $a,b,c$ 和三个夹角 $\alpha,\beta,\gamma$。晶胞体积为
\begin{equation}
V=abc\sqrt{1-\cos^2\alpha-\cos^2\beta-\cos^2\gamma+2\cos\alpha\cos\beta\cos\gamma}.
\label{eq:volume}
\end{equation}
倒易点阵基矢 $\bm{a}^*,\bm{b}^*,\bm{c}^*$ 由实空间基矢 $\bm{a},\bm{b},\bm{c}$ 定义：
\begin{equation}
\bm{a}^*=\frac{\bm{b}\times\bm{c}}{\bm{a}\cdot(\bm{b}\times\bm{c})},\qquad
\bm{b}^*=\frac{\bm{c}\times\bm{a}}{\bm{a}\cdot(\bm{b}\times\bm{c})},\qquad
\bm{c}^*=\frac{\bm{a}\times\bm{b}}{\bm{a}\cdot(\bm{b}\times\bm{c})}.
\label{eq:reciprocal-vectors}
\end{equation}
在晶体学中通常不写固体物理中的 $2\pi$ 因子，因此 $|\bm{a}^*|=1/d_{100}$，即倒易基矢的长度等于同名晶面间距的倒数。倒易点阵的六个参数可用实空间参数表示为
\begin{equation}
a^*=\frac{bc\sin\alpha}{V},\qquad
b^*=\frac{ac\sin\beta}{V},\qquad
c^*=\frac{ab\sin\gamma}{V},
\label{eq:reciprocal-lengths}
\end{equation}
\begin{equation}
\cos\alpha^*=\frac{\cos\beta\cos\gamma-\cos\alpha}{\sin\beta\sin\gamma},\quad
\cos\beta^*=\frac{\cos\alpha\cos\gamma-\cos\beta}{\sin\alpha\sin\gamma},\quad
\cos\gamma^*=\frac{\cos\alpha\cos\beta-\cos\gamma}{\sin\alpha\sin\beta}.
\label{eq:reciprocal-angles}
\end{equation}
对立方晶系，$a=b=c$、$\alpha=\beta=\gamma=90^\circ$，于是
\begin{equation}
V=a^3,\qquad a^*=b^*=c^*=\frac{1}{a},\qquad \alpha^*=\beta^*=\gamma^*=90^\circ.
\label{eq:cubic-reciprocal}
\end{equation}

\subsection{晶面间距、倒易矢量长度与晶面夹角}

晶面 $(hkl)$ 的法向平行于同名倒易矢量 $\bm{G}_{hkl}=h\bm{a}^*+k\bm{b}^*+l\bm{c}^*$，晶面间距等于该倒易矢量长度的倒数：
\begin{equation}
\frac{1}{d_{hkl}^2}=h^2a^{*2}+k^2b^{*2}+l^2c^{*2}
+2hka^*b^*\cos\gamma^*+2hla^*c^*\cos\beta^*+2klb^*c^*\cos\alpha^*.
\label{eq:dspacing}
\end{equation}
对立方晶系，式 \eqref{eq:dspacing} 简化为
\begin{equation}
d_{hkl}=\frac{a}{\sqrt{h^2+k^2+l^2}}.
\label{eq:dspacing-cubic}
\end{equation}
两个倒易矢量 $\bm{G}_1=(h_1k_1l_1)^*$、$\bm{G}_2=(h_2k_2l_2)^*$ 的长度为
\begin{equation}
r_i=\left|\bm{G}_i\right|=\sqrt{h_i^2a^{*2}+k_i^2b^{*2}+l_i^2c^{*2}
+2h_ik_ia^*b^*\cos\gamma^*+2h_il_ia^*c^*\cos\beta^*+2k_il_ib^*c^*\cos\alpha^*},
\quad i=1,2.
\label{eq:rlength}
\end{equation}
二者之间的夹角 $\theta$ 满足
\begin{equation}
\cos\theta=\frac{\bm{G}_1\cdot\bm{G}_2}{r_1r_2}
=\frac{1}{r_1r_2}\left[
\begin{aligned}
&h_1h_2a^{*2}+k_1k_2b^{*2}+l_1l_2c^{*2}\\
&+(k_1l_2+k_2l_1)b^*c^*\cos\alpha^*
 +(l_1h_2+l_2h_1)c^*a^*\cos\beta^*\\
&+(h_1k_2+h_2k_1)a^*b^*\cos\gamma^*
\end{aligned}
\right].
\label{eq:angle}
\end{equation}
式 \eqref{eq:rlength} 和式 \eqref{eq:angle} 就是 Origin 计算模板中 $r_1$、$r_2$ 和 $\theta$ 三列的计算依据。

\subsection{晶带方程与带轴电子衍射}

相交于同一晶向的各晶面族构成一个晶带，该晶向称为晶带轴（zone axis），简称带轴，记为 $[uvw]$。晶面 $(hkl)$ 属于该晶带的条件是
\begin{equation}
hu+kv+lw=0.
\label{eq:zone-law}
\end{equation}
在倒易空间中，属于同一带轴的所有倒易点构成一个过原点的二维倒易面 $(uvw)^*$，带轴 $[uvw]$ 垂直于该面。带轴电子衍射花样就是三维倒易点阵沿该带轴方向的二维截面。由于电子波长很短（200~kV 时约 0.0251~\AA），Ewald 球半径很大，在零阶劳厄区（ZOLZ）内可以近似把倒易面看成与 Ewald 球相切的平面，因此能够同时观察到属于该带轴的大量衍射点。

若在二维倒易面内选取两条不共线的基矢 $\bm{g}_1=(h_1k_1l_1)^*$、$\bm{g}_2=(h_2k_2l_2)^*$，它们的长度分别为 $r_1$、$r_2$，夹角为 $\theta$，则任意格点可用整数 $m,n$ 表示为
\begin{equation}
\bm{g}=m\bm{g}_1+n\bm{g}_2.
\label{eq:2d-lattice}
\end{equation}
以 $\bm{g}_1$ 为 $x$ 轴，则该格点的二维坐标为
\begin{equation}
x=mr_1+nr_2\cos\theta,\qquad y=nr_2\sin\theta,
\label{eq:xy}
\end{equation}
对应的晶面指数为
\begin{equation}
h=mh_1+nh_2,\qquad k=mk_1+nk_2,\qquad l=ml_1+nl_2.
\label{eq:hkl-mn}
\end{equation}
式 \eqref{eq:xy} 和式 \eqref{eq:hkl-mn} 是第 4 章建立衍射图表格的核心。图 \ref{fig:reciprocal-plane} 给出了二维倒易面、基矢和带轴之间的几何关系。

\begin{figure}[H]
  \centering
  \begin{tikzpicture}[scale=1.15,>=stealth]
    \coordinate (O) at (0,0);
    \coordinate (G1) at (4.2,0);
    \coordinate (G2) at (1.6,2.45);
    \coordinate (G12) at (5.8,2.45);
    \draw[->,very thick,secblue] (O)--(G1) node[midway,below=3pt] {$\bm{g}_1$};
    \draw[->,very thick,secblue] (O)--(G2) node[midway,left=3pt] {$\bm{g}_2$};
    \draw[dashed,gray] (G1)--(G12)--(G2);
    \draw[->,very thick,red!70!black] (2.9,1.22)--(2.9,4.15)
      node[above] {$[uvw]\parallel(uvw)^*$};
    \fill (O) circle (1.7pt) node[below left] {000};
    \fill (G1) circle (1.7pt);
    \fill (G2) circle (1.7pt);
    \fill (G12) circle (1.7pt);
    \draw (0.8,0) arc (0:56.9:0.8) node[midway,right] {$\theta$};
    \node[gray] at (2.9,-0.8) {二维倒易面 $(uvw)^*$};
  \end{tikzpicture}
  \caption{带轴电子衍射的二维倒易面、基矢与带轴关系示意。$\bm{g}_1$、$\bm{g}_2$ 为面内基矢，$[uvw]$ 垂直于该倒易面；整数 $m,n$ 将基矢扩展到整个倒易面。}
  \label{fig:reciprocal-plane}
\end{figure}

\subsection{结构因子、消光规律与单晶硅结构}

衍射点的位置由倒易点阵决定，衍射点的强弱则由结构因子决定。晶胞内第 $j$ 个原子对 $(hkl)$ 反射的结构因子贡献为
\begin{equation}
F(hkl)=\sum_{j=1}^{N}g_jt_j(s)f_j(s)\exp\left[2\pi i(hx_j+ky_j+lz_j)\right],
\label{eq:structure-factor}
\end{equation}
其中 $g_j$ 为占有率，$t_j(s)$ 为温度因子，$f_j(s)$ 为原子散射因子，$(x_j,y_j,z_j)$ 为原子分数坐标。原子散射因子描述一个原子在给定方向上的散射振幅相对于一个自由电子散射振幅的比值（量纲为 1），随散射角增大而衰减；在 $s=\sin\theta/\lambda\to0$ 的前向散射极限下，$f(0)$ 近似等于该原子的核外电子数，即原子序数 $Z$。若近似取 $g_j=1$、$t_j=1$、$f_j=f$（常数），则
\begin{equation}
F(hkl)=f\sum_{j=1}^{N}\exp\left[2\pi i(hx_j+ky_j+lz_j)\right],\qquad I\propto|F|^2.
\label{eq:structure-factor-simple}
\end{equation}
本作业的 Origin 模板正是采用这一简化形式，并取 Si 的 $f=14$（Si 的原子序数 $Z=14$，即取前向散射极限值并忽略其随角度的衰减，详见第 \ref{sec:f14} 节）。

单晶硅为金刚石型立方结构，空间群 $Fd\bar{3}m$（No.~227），点阵常数 $a=5.43$~\AA，每个惯用晶胞含 8 个 Si 原子，其分数坐标列于表 \ref{tab:si-atoms}。该结构可看作两套面心立方点阵沿体对角线相对位移 $(1/4,1/4,1/4)$ 套构而成，每个 Si 原子与 4 个近邻形成正四面体配位，Si--Si 键长为 $\sqrt{3}a/4\approx2.35$~\AA。

\begin{table}[H]
  \centering
  \caption{单晶硅（$Fd\bar{3}m$，$a=5.43$~\AA）晶胞内 8 个 Si 原子的分数坐标}
  \label{tab:si-atoms}
  \small
  \begin{tabular}{cccc}
    \toprule
    序号 & $x$ & $y$ & $z$ \\
    \midrule
    1 & 0     & 0     & 0     \\
    2 & 0     & 1/2   & 1/2   \\
    3 & 1/2   & 0     & 1/2   \\
    4 & 1/2   & 1/2   & 0     \\
    5 & 1/4   & 1/4   & 1/4   \\
    6 & 1/4   & 3/4   & 3/4   \\
    7 & 3/4   & 1/4   & 3/4   \\
    8 & 3/4   & 3/4   & 1/4   \\
    \bottomrule
  \end{tabular}
\end{table}

把表 \ref{tab:si-atoms} 的坐标代入式 \eqref{eq:structure-factor-simple}，可得金刚石结构的结构因子
\begin{equation}
F=f\left[1+e^{\pi i(h+k)}+e^{\pi i(h+l)}+e^{\pi i(k+l)}\right]
\left[1+e^{\frac{\pi i}{2}(h+k+l)}\right].
\label{eq:diamond-F}
\end{equation}
由此得到金刚石结构的消光规律：
\begin{itemize}[leftmargin=2em,itemsep=1pt]
  \item $h,k,l$ 奇偶混合时，第一个方括号为 0，$F=0$；
  \item $h,k,l$ 全为偶数且 $h+k+l=4n+2$ 时，第二个方括号为 0，$F=0$；
  \item $h,k,l$ 全为奇数时，$|F|^2=32f^2$；
  \item $h,k,l$ 全为偶数且 $h+k+l=4n$ 时，$|F|^2=64f^2$。
\end{itemize}
因此，Si 的允许反射为 111、220、311、400、331、422、511/333、440 等，而 100、110、200、222、420、600 等消光。第 3 章将以 110 反射为例，用模板计算验证这一规律。
